% =============================================================================
% APPENDIX SG: REF-STYLE: FehrAdvice Corporate Style Guide
% =============================================================================
% Module: BCM2_00_META
% Version: 1.0 (January 2026)
% Status: Final
%
% Single Source of Truth für Corporate Identity und Design-Standards
% von FehrAdvice \& Partners AG. Definiert Farben, Typografie, Layout
% und Kommunikationsrichtlinien für alle EBF-Outputs.
%
% =============================================================================

% =============================================================================
% CROSS-REFERENCE STRUCTURE
% =============================================================================
%
%                      ┌─────────────────────┐
%                      │   Appendix SG       │
%                      │   REF-STYLE         │
%                      │   Style Guide       │
%                      └──────────┬──────────┘
%                                 │
%           ┌─────────────────────┼─────────────────────┐
%           │                     │                     │
%           ▼                     ▼                     ▼
%   ┌───────────────┐    ┌───────────────┐    ┌───────────────┐
%   │ Appendix CCC  │    │ Appendix DDD  │    │ Appendix G    │
%   │ METHOD-DOCTYPE│    │ REF-DOCTYPE   │    │ REF-GLOSSARY  │
%   │ (8D-Algorithm)│    │ (Doc Types)   │    │ (Terminology) │
%   └───────────────┘    └───────────────┘    └───────────────┘
%
% PRIMARY CHAPTER: None (Meta-Reference)
% SECONDARY CHAPTERS: All (applies to all outputs)
%
% =============================================================================

\chapter{FehrAdvice Corporate Style Guide}
\label{app:ref-sg}

% -----------------------------------------------------------------------------
% ABSTRACT
% -----------------------------------------------------------------------------

\begin{abstract}
Dieser Appendix definiert die verbindlichen Corporate-Identity-Standards für
FehrAdvice \& Partners AG. Er enthält die offizielle Farbpalette, Typografie-
spezifikationen, Layout-Richtlinien und Kommunikationsstandards für alle
EBF-Outputs, Reports und Präsentationen. Als Single Source of Truth (SSOT)
stellt er konsistente Markenidentität über alle Dokumente sicher.
\end{abstract}

% -----------------------------------------------------------------------------
% QUICK REFERENCE BOX
% -----------------------------------------------------------------------------

\begin{tcolorbox}[colback=blue!5!white, colframe=blue!75!black, title=Quick Reference: FehrAdvice Farben]
\begin{tabular}{llll}
\textbf{Primärfarben} & \textbf{HEX} & \textbf{RGB} & \textbf{Verwendung} \\
\hline
Dunkelblau & \texttt{\#024079} & RGB(2,64,121) & Headlines, Links, Akzente \\
Hellblau & \texttt{\#549EDE} & RGB(84,158,222) & Sekundäre Elemente \\
Dunkelgrau & \texttt{\#25212A} & RGB(37,33,42) & Fliesstext \\
Hellgrau & \texttt{\#F3F5F7} & RGB(243,245,247) & Hintergründe \\
\hline
\textbf{Sekundärfarben} & & & \\
\hline
Flieder & \texttt{\#A1A0C6} & RGB(161,160,198) & Akzente, Diagramme \\
Mintgrün & \texttt{\#7EBDAC} & RGB(126,189,172) & Akzente, Diagramme \\
Ocker & \texttt{\#DECB3F} & RGB(222,203,63) & Hervorhebungen \\
Orange & \texttt{\#DE9D3E} & RGB(222,157,62) & Warnungen, CTA \\
\end{tabular}
\end{tcolorbox}

% -----------------------------------------------------------------------------
% APPENDIX SCOPE BOX
% -----------------------------------------------------------------------------

\begin{tcolorbox}[colback=gray!5!white, colframe=gray!75!black, title=Appendix Scope]
\textbf{Ziel:} Verbindliche Definition aller visuellen und sprachlichen Standards für FehrAdvice-Outputs.

\textbf{In-Scope:}
\begin{itemize}
    \item Farbpalette (Primär- und Sekundärfarben mit HEX/RGB)
    \item Typografie (Schriftarten, -grössen, -gewichte)
    \item Layout-Spezifikationen (Seitenformat, Ränder, Abstände)
    \item Tabellengestaltung und Diagrammfarben
    \item Logo-Verwendung und Schutzraum
    \item Tone of Voice und Schreibstil
    \item Schweizer Orthographie und Gendering
\end{itemize}

\textbf{Out-of-Scope:}
\begin{itemize}
    \item Dokumenttyp-Auswahl (→ Appendix CCC/DDD: 8D-Algorithmus)
    \item Content-Struktur (→ Appendix Templates)
    \item Wissenschaftliche Zitierweise (→ bcm\_master.bib)
\end{itemize}

\textbf{Lieferobjekte:}
\begin{itemize}
    \item LaTeX-Farbdefinitionen für alle Templates
    \item CSS-Spezifikationen für Web/PDF-Outputs
    \item Checkliste für Style-Compliance
\end{itemize}
\end{tcolorbox}


% =============================================================================
%                           PART 2: CORE CONTENT
% =============================================================================

\section{Farbpalette}
\label{sec:sg-colors}

Die FehrAdvice-Farbpalette basiert auf einer professionellen, vertrauenswürdigen
Ästhetik, die wissenschaftliche Seriosität mit moderner Zugänglichkeit verbindet.

\subsection{Primärfarben}

\begin{table}[h]
\centering
\caption{FehrAdvice Primärfarben}
\label{tab:primary-colors}
\begin{tabular}{llll}
\toprule
\textbf{Farbe} & \textbf{HEX} & \textbf{RGB} & \textbf{Verwendung} \\
\midrule
Dunkelblau & \texttt{\#024079} & RGB(2,64,121) & Überschriften, Links, Akzente \\
Hellblau & \texttt{\#549EDE} & RGB(84,158,222) & Sekundäre Elemente \\
Dunkelgrau & \texttt{\#25212A} & RGB(37,33,42) & Fliesstext \\
Hellgrau & \texttt{\#F3F5F7} & RGB(243,245,247) & Hintergründe \\
\bottomrule
\end{tabular}
\end{table}

\subsection{Sekundärfarben}

\begin{table}[h]
\centering
\caption{FehrAdvice Sekundärfarben}
\label{tab:secondary-colors}
\begin{tabular}{llll}
\toprule
\textbf{Farbe} & \textbf{HEX} & \textbf{RGB} & \textbf{Verwendung} \\
\midrule
Flieder & \texttt{\#A1A0C6} & RGB(161,160,198) & Akzente, Diagramme \\
Mintgrün & \texttt{\#7EBDAC} & RGB(126,189,172) & Akzente, Diagramme \\
Ocker & \texttt{\#DECB3F} & RGB(222,203,63) & Hervorhebungen \\
Orange & \texttt{\#DE9D3E} & RGB(222,157,62) & Warnungen, Call-to-Action \\
\bottomrule
\end{tabular}
\end{table}

\subsection{Diagrammreihenfolge}

Für Diagramme und Charts mit mehreren Datenreihen gilt folgende Farbreihenfolge:

\begin{enumerate}
    \item Dunkelblau (\texttt{\#024079})
    \item Hellblau (\texttt{\#549EDE})
    \item Mintgrün (\texttt{\#7EBDAC})
    \item Flieder (\texttt{\#A1A0C6})
    \item Ocker (\texttt{\#DECB3F})
    \item Orange (\texttt{\#DE9D3E})
\end{enumerate}

\subsection{LaTeX-Farbdefinitionen}

\begin{verbatim}
% Primärfarben
\definecolor{primarydarkblue}{RGB}{2,64,121}
\definecolor{primarylightblue}{RGB}{84,158,222}
\definecolor{primarydarkgray}{RGB}{37,33,42}
\definecolor{primarylightgray}{RGB}{243,245,247}

% Sekundärfarben
\definecolor{secondarylilac}{RGB}{161,160,198}
\definecolor{secondarymint}{RGB}{126,189,172}
\definecolor{secondaryocher}{RGB}{222,203,63}
\definecolor{secondaryorange}{RGB}{222,157,62}
\end{verbatim}


% =============================================================================
\section{Typografie}
\label{sec:sg-typography}

\subsection{Schriftarten}

\begin{table}[h]
\centering
\caption{FehrAdvice Schriftarten}
\label{tab:fonts}
\begin{tabular}{lll}
\toprule
\textbf{Verwendung} & \textbf{Schriftart} & \textbf{Fallback} \\
\midrule
Headlines (H1, H2) & Roboto Bold & Arial Bold, sans-serif \\
Subheadlines (H3, H4) & Roboto Regular & Arial, sans-serif \\
Fliesstext & Open Sans Regular & Arial, sans-serif \\
Tabellen \& Daten & Open Sans Regular & Arial, sans-serif \\
Akzente \& Zitate & Playfair Display & Georgia, serif \\
Code \& Technisch & Consolas & Monaco, monospace \\
\bottomrule
\end{tabular}
\end{table}

\textbf{Schriftarten-Download:}
\begin{itemize}
    \item Roboto: \texttt{fonts.google.com/specimen/Roboto}
    \item Open Sans: \texttt{fonts.google.com/specimen/Open+Sans}
    \item Playfair Display: \texttt{fonts.google.com/specimen/Playfair+Display}
\end{itemize}

\subsection{Schriftgrössen}

\begin{table}[h]
\centering
\caption{FehrAdvice Schriftgrössen}
\label{tab:font-sizes}
\begin{tabular}{llllll}
\toprule
\textbf{Element} & \textbf{Font} & \textbf{Grösse} & \textbf{Zeilenabstand} & \textbf{Gewicht} & \textbf{Farbe} \\
\midrule
H1 – Haupttitel & Roboto & 28 pt & 34 pt & Bold & Dunkelblau \\
H2 – Kapitelüberschrift & Roboto & 22 pt & 28 pt & Bold & Dunkelblau \\
H3 – Abschnittstitel & Roboto & 16 pt & 22 pt & Regular & Dunkelblau \\
H4 – Untertitel & Roboto & 13 pt & 18 pt & Regular & Dunkelblau \\
Body – Fliesstext & Open Sans & 11 pt & 16 pt & Regular & Dunkelgrau \\
Caption – Bildunterschriften & Open Sans & 9 pt & 12 pt & Regular & Dunkelgrau \\
Fussnoten & Open Sans & 8 pt & 10 pt & Regular & Dunkelgrau \\
Akzente/Zitate & Playfair Display & variabel & variabel & Regular & Dunkelblau \\
\bottomrule
\end{tabular}
\end{table}


% =============================================================================
\section{Layout \& Raster}
\label{sec:sg-layout}

\subsection{Seitenformat}

\begin{table}[h]
\centering
\caption{Seitenformat-Spezifikationen}
\label{tab:page-format}
\begin{tabular}{ll}
\toprule
\textbf{Element} & \textbf{Wert} \\
\midrule
Format & A4 (210 $\times$ 297 mm) \\
Ausrichtung & Hochformat (Portrait) \\
Seitenrand oben & 25 mm \\
Seitenrand unten & 20 mm \\
Seitenrand links & 25 mm \\
Seitenrand rechts & 20 mm \\
\bottomrule
\end{tabular}
\end{table}

\subsection{Abstände}

\begin{table}[h]
\centering
\caption{Abstands-Spezifikationen}
\label{tab:spacing}
\begin{tabular}{ll}
\toprule
\textbf{Element} & \textbf{Abstand} \\
\midrule
Nach H1 & 18 pt \\
Nach H2 & 14 pt \\
Nach H3 & 10 pt \\
Zwischen Absätzen & 10 pt \\
Vor Listen & 6 pt \\
Listeneinzug & 15 pt \\
\bottomrule
\end{tabular}
\end{table}


% =============================================================================
\section{Visuelle Elemente}
\label{sec:sg-visual}

\subsection{Tabellen}

\begin{table}[h]
\centering
\caption{Tabellen-Spezifikationen}
\label{tab:table-specs}
\begin{tabular}{ll}
\toprule
\textbf{Element} & \textbf{Spezifikation} \\
\midrule
Kopfzeile Hintergrund & Dunkelblau (\texttt{\#024079}) \\
Kopfzeile Schrift & Weiss, Bold, 9 pt \\
Zeilen alternierend & Weiss / Hellgrau (\texttt{\#F3F5F7}) \\
Rahmenlinien & Hellgrau (\texttt{\#E0E0E0}), 0.5 pt \\
Zellenpadding & 8 pt \\
\bottomrule
\end{tabular}
\end{table}

\subsection{Diagramme \& Charts}

\begin{table}[h]
\centering
\caption{Diagramm-Spezifikationen}
\label{tab:chart-specs}
\begin{tabular}{ll}
\toprule
\textbf{Element} & \textbf{Spezifikation} \\
\midrule
Primäre Datenreihe & Dunkelblau (\texttt{\#024079}) \\
Sekundäre Datenreihen & Hellblau, Mintgrün, Flieder, Ocker, Orange \\
Achsenbeschriftung & Dunkelgrau, 9 pt \\
Gitternetzlinien & Hellgrau (\texttt{\#E0E0E0}), 0.5 pt \\
Legende & Rechts oder unten, 9 pt \\
\bottomrule
\end{tabular}
\end{table}

\subsection{Icons}

\begin{itemize}
    \item \textbf{Stil:} Liniensymbole (Outline), nicht gefüllt
    \item \textbf{Strichstärke:} 1.5–2 pt
    \item \textbf{Farbe:} Dunkelblau oder Dunkelgrau
    \item \textbf{Grösse:} 24 $\times$ 24 px (Standard)
\end{itemize}


% =============================================================================
\section{Logo}
\label{sec:sg-logo}

\subsection{Geschichte und Bedeutung: Der Schlussstein}

Das Logo von FehrAdvice, symbolisiert durch den \textbf{Schlussstein}, fungiert
als authentischer Ausdruck der Unternehmenswerte. Der Schlussstein, in der
Architektur bekannt für seine Fähigkeit, Bögen zusammenzuhalten, symbolisiert
bei FehrAdvice die Kollaboration und den kollektiven Fortschritt.

Das Logo reflektiert auch die Brückenfunktion von FehrAdvice zwischen
akademischer Forschung und praktischer Anwendung.

\subsection{Logo-Aufbau}

\begin{table}[h]
\centering
\caption{Logo-Elemente}
\label{tab:logo-elements}
\begin{tabular}{ll}
\toprule
\textbf{Element} & \textbf{Beschreibung} \\
\midrule
Symbol & Der Schlussstein – stilisierte, gestaffelte Bögen \\
Wortmarke & «FEHR ADVICE» in Versalien \\
Tagline & «Behavioral Economics Consultancy Group» \\
\bottomrule
\end{tabular}
\end{table}

\subsection{Logo-Farben}

\begin{table}[h]
\centering
\caption{Logo-Farbspezifikationen}
\label{tab:logo-colors}
\begin{tabular}{llll}
\toprule
\textbf{Element} & \textbf{Farbe} & \textbf{HEX} & \textbf{RGB} \\
\midrule
Symbol (Schlussstein) & Dunkelblau & \texttt{\#024079} & RGB(2,64,121) \\
Wortmarke & Dunkelblau & \texttt{\#024079} & RGB(2,64,121) \\
Tagline & Dunkelblau & \texttt{\#024079} & RGB(2,64,121) \\
Monochrom Weiss & Weiss & \texttt{\#FFFFFF} & RGB(255,255,255) \\
\bottomrule
\end{tabular}
\end{table}

\subsection{Verwendung mit/ohne Tagline}

\begin{table}[h]
\centering
\caption{Tagline-Verwendung}
\label{tab:tagline-usage}
\begin{tabular}{ll}
\toprule
\textbf{Kontext} & \textbf{Empfehlung} \\
\midrule
Mit Tagline & Offizielle Dokumente, Präsentationen, Berichte, Website \\
Ohne Tagline & Kleine Anwendungen, Favicon, Social Media, interne Docs \\
\bottomrule
\end{tabular}
\end{table}

\subsection{Mindestgrössen und Schutzraum}

\begin{itemize}
    \item \textbf{Print:} Mindestbreite 25 mm
    \item \textbf{Digital:} Mindestbreite 100 px
    \item \textbf{Favicon:} 16 $\times$ 16 px
    \item \textbf{Schutzraum:} Mindestens Höhe des «F» in der Wortmarke
\end{itemize}

\subsection{Logo Don'ts}

\begin{itemize}
    \item Logo nicht strecken oder verzerren
    \item Logo nicht drehen oder kippen
    \item Logo nicht in nicht-definierten Farben darstellen
    \item Logo nicht auf unruhigen Hintergründen platzieren
    \item Logo nicht mit Schatten oder Effekten versehen
    \item Mindestabstand immer einhalten
\end{itemize}


% =============================================================================
\section{Tone of Voice}
\label{sec:sg-tone}

\subsection{Markenpersönlichkeit}

FehrAdvice kommuniziert:

\begin{table}[h]
\centering
\caption{Markenpersönlichkeit}
\label{tab:brand-personality}
\begin{tabular}{ll}
\toprule
\textbf{Eigenschaft} & \textbf{Beschreibung} \\
\midrule
Wissenschaftlich fundiert & Aussagen basieren auf Forschung und Evidenz \\
Klar und präzise & Komplexe Themen verständlich erklärt \\
Partnerschaftlich & Auf Augenhöhe mit Kund:innen \\
Zukunftsorientiert & Fokus auf nachhaltige Lösungen \\
Schweizerisch & Qualität, Zuverlässigkeit, Diskretion \\
\bottomrule
\end{tabular}
\end{table}

\subsection{Schreibstil}

\begin{table}[h]
\centering
\caption{Schreibstil Do's und Don'ts}
\label{tab:writing-style}
\begin{tabular}{ll}
\toprule
\textbf{Do} & \textbf{Don't} \\
\midrule
Aktive Formulierungen & Passive Konstruktionen \\
Kurze, klare Sätze & Verschachtelte Sätze \\
Konkrete Beispiele & Abstrakte Theorie ohne Bezug \\
Evidenzbasierte Aussagen & Unbelegte Behauptungen \\
«Wir empfehlen...» & «Man sollte...» \\
\bottomrule
\end{tabular}
\end{table}

\subsection{Kernbegriffe (Do's)}

\begin{itemize}
    \item Verhaltensökonomie
    \item Behavioral Insights
    \item Evidenzbasiert
    \item Nudging
    \item Decision Architecture
    \item Behavioral Design
    \item Customer Centricity
    \item Human-Centered
\end{itemize}

\subsection{Begriffe vermeiden (Don'ts)}

\begin{itemize}
    \item Manipulation
    \item Trick
    \item Austricksen
    \item Überlisten
    \item Zwingen
    \item Müssen (wenn möglich durch «empfehlen» ersetzen)
\end{itemize}


% =============================================================================
\section{Orthographie}
\label{sec:sg-orthography}

\subsection{Schweizer Schreibweise}

FehrAdvice \& Partners AG verwendet die schweizerische Orthographie:

\begin{table}[h]
\centering
\caption{Schweizer vs. Deutsche Schreibweise}
\label{tab:swiss-orthography}
\begin{tabular}{lll}
\toprule
\textbf{Element} & \textbf{Schweiz} & \textbf{Deutschland} \\
\midrule
Anführungszeichen & « » & „ " \\
Doppel-S & ss & ß \\
Tausendertrennzeichen & 1'000.00 & 1.000,00 \\
\bottomrule
\end{tabular}
\end{table}

\textbf{Beispiele:}
\begin{itemize}
    \item «Das ist ein Zitat.»
    \item Die Strasse führt zum Fluss.
    \item CHF 1'250'000.00
\end{itemize}

\subsection{Gendering}

FehrAdvice verwendet den Doppelpunkt als Gender-Separator:

\begin{itemize}
    \item Kund:innen
    \item Bürger:innen
    \item Mitarbeiter:innen
    \item Berater:innen
\end{itemize}

\subsection{Englische Texte}

FehrAdvice verwendet \textbf{amerikanisches Englisch} (American English):

\begin{table}[h]
\centering
\caption{US vs. UK Englisch}
\label{tab:us-english}
\begin{tabular}{ll}
\toprule
\textbf{Korrekt (US)} & \textbf{Falsch (UK)} \\
\midrule
Color & Colour \\
Organize & Organise \\
Behavior & Behaviour \\
Center & Centre \\
Analyze & Analyse \\
\bottomrule
\end{tabular}
\end{table}

\subsection{Headlines (Englisch)}

FehrAdvice verwendet die \textbf{Kleinschreibvariante} für englische Headlines:

\textbf{Korrekt:} \textit{This is a headline where the default rule is lower case except for America}

\textbf{Falsch:} \textit{This Is a Headline With Complicated Spelling of Most First Letters}

\textbf{Regel:} Nur Satzanfänge und Eigennamen werden gross geschrieben.


% =============================================================================
\section{Barrierefreiheit}
\label{sec:sg-accessibility}

\subsection{Farbkontrast}

\begin{table}[h]
\centering
\caption{Farbkontrast-Prüfung (WCAG)}
\label{tab:color-contrast}
\begin{tabular}{lll}
\toprule
\textbf{Kombination} & \textbf{Kontrastverhältnis} & \textbf{Status} \\
\midrule
Dunkelblau auf Weiss & 10.5:1 & $\checkmark$ AAA \\
Dunkelgrau auf Weiss & 12.6:1 & $\checkmark$ AAA \\
Weiss auf Dunkelblau & 10.5:1 & $\checkmark$ AAA \\
Dunkelblau auf Hellgrau & 9.8:1 & $\checkmark$ AAA \\
\bottomrule
\end{tabular}
\end{table}

\subsection{Mindestgrössen}

\begin{itemize}
    \item \textbf{Fliesstext:} Minimum 11 pt
    \item \textbf{Fussnoten:} Minimum 8 pt
    \item \textbf{Buttons/Links (klickbar):} Minimum 44 $\times$ 44 px
\end{itemize}

\subsection{Alternativtexte}

\begin{itemize}
    \item Alle Bilder benötigen beschreibende Alt-Texte
    \item Diagramme benötigen Textbeschreibungen der Kernaussage
    \item Dekorative Elemente: leerer Alt-Text (\texttt{alt=""})
\end{itemize}


% =============================================================================
%                           PART 3: APPLICATION
% =============================================================================

\section{Worked Example: PDF-Report Styling}
\label{sec:sg-example}

\subsection{CSS für WeasyPrint-Outputs}

Das folgende CSS implementiert den FehrAdvice Style Guide für PDF-Generierung:

\begin{verbatim}
@import url('https://fonts.googleapis.com/css2?family=Open+Sans&family=Playfair+Display&family=Roboto:wght@400;700&display=swap');

@page {
    size: A4;
    margin: 25mm 20mm 20mm 25mm;
}

body {
    font-family: "Open Sans", Arial, sans-serif;
    font-size: 11pt;
    line-height: 16pt;
    color: #25212A;
}

h1, h2, h3, h4 {
    font-family: "Roboto", Arial, sans-serif;
    color: #024079;
}

h1 {
    font-size: 28pt;
    font-weight: 700;
    border-bottom: 2px solid #024079;
}

h2 {
    font-size: 22pt;
    font-weight: 700;
    border-bottom: 1px solid #549EDE;
}

blockquote {
    font-family: "Playfair Display", Georgia, serif;
    font-style: italic;
    color: #024079;
}

th {
    background-color: #024079;
    color: #FFFFFF;
    font-family: "Roboto", Arial, sans-serif;
}

tr:nth-child(even) {
    background-color: #F3F5F7;
}
\end{verbatim}

\subsection{Checkliste für Style-Compliance}

\begin{tcolorbox}[colback=green!5!white, colframe=green!75!black, title=Style-Compliance Checkliste]
\begin{itemize}
    \item[$\square$] Headlines in \textbf{Roboto Bold}, Dunkelblau (\texttt{\#024079})
    \item[$\square$] Fliesstext in \textbf{Open Sans}, Dunkelgrau (\texttt{\#25212A})
    \item[$\square$] Akzente/Zitate in \textbf{Playfair Display} (optional)
    \item[$\square$] Tabellenkopf Dunkelblau mit weisser Schrift
    \item[$\square$] Alternating rows: Weiss/Hellgrau (\texttt{\#F3F5F7})
    \item[$\square$] Seitenränder: 25/20/25/20 mm
    \item[$\square$] Schweizer Anführungszeichen « »
    \item[$\square$] Gendering mit Doppelpunkt (:)
    \item[$\square$] Amerikanisches Englisch (behavior, nicht behaviour)
    \item[$\square$] Logo mit korrektem Schutzraum
\end{itemize}
\end{tcolorbox}


% =============================================================================
%                           PART 4: BACK MATTER
% =============================================================================

\section{Summary}
\label{sec:sg-summary}

\begin{tcolorbox}[colback=blue!5!white, colframe=blue!75!black, title=Summary: FehrAdvice Style Guide]
\textbf{Kernelemente:}
\begin{itemize}
    \item \textbf{Primärfarbe:} Dunkelblau \texttt{\#024079} für alle Headlines und Akzente
    \item \textbf{Textfarbe:} Dunkelgrau \texttt{\#25212A} für Fliesstext
    \item \textbf{Hintergrund:} Hellgrau \texttt{\#F3F5F7} für Tabellen-Alternation
    \item \textbf{Fonts:} Roboto (Headlines), Open Sans (Body), Playfair Display (Akzente)
    \item \textbf{Layout:} A4, Ränder 25/20/25/20 mm
    \item \textbf{Sprache:} Schweizer Orthographie, US-Englisch, Gender-Doppelpunkt
    \item \textbf{Tone:} Wissenschaftlich fundiert, klar, partnerschaftlich
\end{itemize}

\textbf{Quelle:} Offizieller Corporate Identity Guide (assets/Corporate-Identity-Guide-FAP-EN-231007 2.pdf)

\textbf{Single Source of Truth:} Dieser Appendix ist verbindlich für alle EBF-Outputs.
\end{tcolorbox}


% -----------------------------------------------------------------------------
% GLOSSARY SECTION
% -----------------------------------------------------------------------------

\section{Glossary of Terms}
\label{sec:sg-glossary}

\begin{tabular}{ll}
\textbf{Term} & \textbf{Definition} \\
\hline
CI/CD & Corporate Identity / Corporate Design \\
HEX & Hexadezimale Farbnotation (\texttt{\#RRGGBB}) \\
RGB & Rot-Grün-Blau Farbmodell \\
WCAG & Web Content Accessibility Guidelines \\
AAA & Höchster WCAG-Kontrast-Standard \\
CTA & Call-to-Action \\
\end{tabular}

\vspace{10pt}
\textit{Für weitere Definitionen siehe Appendix G (Master Glossary).}


% -----------------------------------------------------------------------------
% REFERENCES
% -----------------------------------------------------------------------------

\section{References}
\label{sec:sg-references}

\nocite{bcm_master}

Dieser Style Guide basiert auf:
\begin{itemize}
    \item \textbf{FehrAdvice Corporate Identity Guide} (October 2023): \texttt{assets/Corporate-Identity-Guide-FAP-EN-231007 2.pdf}
    \item Web Content Accessibility Guidelines (WCAG) 2.1
    \item Swiss Typography Standards (Schweizer Typographische Gesellschaft)
    \item Google Fonts (Roboto, Open Sans, Playfair Display)
\end{itemize}


% -----------------------------------------------------------------------------
% CROSS-REFERENCE FOOTER
% -----------------------------------------------------------------------------

\vspace{20pt}
\hrule
\vspace{10pt}

\textbf{Cross-References:}
\begin{itemize}
    \item \textbf{Appendix CCC (METHOD-DOCTYPE):} 8D-Algorithmus für Dokumenttyp-Auswahl
    \item \textbf{Appendix DDD (REF-DOCTYPE):} Dokumenttyp-Definitionen
    \item \textbf{Appendix G (REF-GLOSSARY):} Master Glossary
\end{itemize}

\vspace{10pt}
\textit{FehrAdvice \& Partners AG – Stand: Januar 2026}

\end{chapter}
